% Ad Atticum 9.10 (ad-atticum-09-10)
% Source: https://raw.githubusercontent.com/PerseusDL/canonical-latinLit/master/data/phi0474/phi057/phi0474.phi057.perseus-lat1.xml
% Fetched by scripts/fetch_latin.py

% Dateline (Perseus): Scr. in Formiano xv K. Apr. a. 705 (49) .

\ciceroLetterOpener{CICERO ATTICO salutem}

\ciceroSection{1} nihil habebam quod scriberem. neque enim novi quicquam audieram et ad tuas omnis rescripseram pridie. sed cum me aegritudo non solum somno privaret verum ne vigilare quidem sine summo dolore pateretur, tecum ut quasi loquerer, in quo uno acquiesco, hoc nescio quid nullo argumento proposito scribere institui.

\ciceroSection{2} amens mihi fuisse videor a principio et me una haec res torquet quod non omnibus in rebus labentem vel potius ruentem Pompeium tamquam unus manipularis secutus sim. vidi hominem xiiii K. Febr. plenum formidinis. illo ipso die sensi quid ageret. numquam mihi postea placuit nec umquam aliud in alio peccare destitit. nihil interim ad me scribere, nihil nisi fugam cogitare. quid quaeris? sicut ἐν τοῖσ ἐρωτικοῖσ alienant immundae, insulsae, indecorae, sic me illius fugae neglegentiaeque deformitas avertit ab amore. nihil enim dignum faciebat qua re eius fugae comitem me adiungerem. nunc emergit amor, nunc desiderium ferre non possum, nunc mihi nihil libri, nihil litterae, nihil doctrina prodest. ita dies et noctes tamquam avis illa mare prospecto evolare cupio. do , do poenas temeritatis meae. etsi quae fuit illa temeritas? quid feci non consideratissime? si enim nihil praeter fugam quaereretur, fugissem libentissime, sed genus belli crudelissimi et maximi quod nondum vident homines quale futurum sit perhorrui. quae minae municipiis, quae nominatim viris bonis, quae denique omnibus qui remansissent! quam crebro illud Sulla potuit, ego non potero?

\ciceroSection{3} mihi autem haeserunt illa. male Tarquinius qui Porsenam, qui Octavium Mamilium contra patriam, impie Coriolanus qui auxilium petiit a Volscis, recte Themistocles qui mori maluit, nefarius Hippias Pisistrati filius qui in Marathonia pugna cecidit arma contra patriam ferens; at Sulla, at Marius, at Cinna recte, immo iure fortasse; sed quid eorum victoria crudelius, quid funestius? huius belli genus fugi et eo magis quod crudeliora etiam cogitari et parari videbam. me quem non nulli conservatorem istius urbis, quem parentem esse dixerunt Getarum et Armeniorum et Colchorum copias ad eam adducere? me meis civibus famem, vastitatem inferre Italiae? hunc primum mortalem esse, deinde etiam multis modis posse exstingui cogitabam, urbem autem et populum nostrum servandum ad immortalitatem, quantum in nobis esset, putabam, et tamen spes quaedam me sustentabat fore ut aliquid conveniret potius quam aut hic tantum sceleris aut ille tantum flagiti admitteret. alia res nunc tota est, alia mens mea. sol , ut est in tua quadam epistula, excidisse mihi e mundo videtur. ut aegroto, dum anima est, spes esse dicitur, sic ego, quoad Pompeius in Italia fuit, sperare non destiti. haec , haec me fefellerunt et, ut verum loquar, aetas iam a diuturnis laboribus devexa ad otium domesticarum me rerum delectatione mollivit. nunc si vel periculose experiundum erit, experiar certe ut hinc avolem. ante oportuit fortasse; sed ea quae scripsisti me tardarunt et auctoritas maxime tua.

\ciceroSection{4} nam cum ad hunc locum venissem, evolvi volumen epistularum tuarum quod ego sub signo habeo servoque diligentissime. erat igitur in ea quam x K. Febr. dederas, hoc modo, sed videamus et Gnaeus quid agat et illius rationes quorsum fluant. quod si iste Italiam relinquet, faciet omnino male et, ut ego existimo, ἀλογίστωσ , sed tum demum consilia nostra commutanda erunt. hoc scribis post diem quartum quam ab urbe discessimus. deinde viii K. Febr. , tantum modo Gnaeus noster ne, ut urbem ἀλογίστωσ reliquit, sic Italiam relinquat. eodem die das alteras litteras quibus mihi consulenti planissime respondes. est enim sic, sed venio ad consultationem tuam. si Gnaeus Italia cedit, in urbem redeundum puto; quae enim finis peregrinationis? hoc mihi plane haesit, et nunc ita video infinitum bellum iunctum miserrima fuga quam tu peregrinationem ὑποκορίζῃ .

\ciceroSection{5} sequitur χρησμὸσ vi K. Februarias, ego si Pompeius manet in Italia nec res ad pactionem venit, longius bellum puto fore; sin Italiam relinquit, ad posterum bellum ἄσπονδον strui existimo. huius igitur belli ego particeps et socius et adiutor esse cogor quod et ἄσπονδον est et cum civibus? deinde vii Idus Febr. , cum iam plura audires de Pompei consilio, concludis epistulam quandam hoc modo, ego quidem tibi non sim auctor, si Pompeius Italiam relinquit, te quoque profugere. summo enim periculo facies nec rei publicae proderis; quoi quidem posterius poteris prodesse, si manseris. quem φιλόπατριν ac πολιτικὸν hominis prudentis et amici tali admonitu non moveret auctoritas?

\ciceroSection{6} deinceps iii Idus Febr. iterum mihi respondes consulenti sic, quod quaeris a me fugamne defendam an moram utiliorem putem, ego vero in praesentia subitum discessum et praecipitem profectionem cum tibi tum ipsi Gnaeo inutilem et periculosam puto et satius esse existimo vos dispertitos et in speculis esse; sed medius fidius turpe nobis puto esse de fuga cogitare. hoc turpe Gnaeus noster biennio ante cogitavit. ita sullaturit animus eius et proscripturit iam diu. inde , ut opinor, cum tu ad me quaedam γενικώτερον scripsisses et ego mihi a te quaedam significari putassem ut Italia cederem, detestaris hoc diligenter xi K. Mart. , ego vero nulla epistula significavi, si Gnaeus Italia cederet, ut tu una cederes, aut si significavi, non dico fui inconstans sed demens. in eadem epistula alio loco, nihil relinquitur nisi fuga, cui te socium neutiquam puto esse oportere nec umquam putavi.

\ciceroSection{7} totam autem hanc deliberationem evolvis accuratius in litteris viii Kal. Mart. datis, si M'. Lepidus et L. Volcacius remanent, manendum puto, ita ut, si salvus sit Pompeius et constiterit alicubi, hanc νέκυιαν relinquas et te in certamine vinci cum illo facilius patiaris quam cum hoc in ea quae perspicitur futura colluvie regnare. multa disputas huic sententiae convenientia. inde ad extremum, quid si inquis Lepidus et Volcacius discedunt? plane ἀπορῶ . quod evenerit igitur et quod egeris id στερκτέον putabo. si tum dubitaras, nunc certe non dubitas istis manentibus.

\ciceroSection{8} deinde in ipsa fuga v Kal. Martias , interea non dubito quin in Formiano mansurus sis. commodissime enim τὸ μέλλον ibi καραδοκήσεισ . atque K. Mart., cum ille quintum iam diem Brundisi esset, tum poterimus deliberare non scilicet integra re sed certe minus infracta quam si una proieceris te. deinde iiii Non. Martias , ὑπὸ τὴν λῆψιν cum breviter scriberes, tamen ponis hoc, cras scribam plura et ad omnia; hoc tamen dicam, non paenitere me consili de tua mansione, et quamquam magna sollicitudine, tamen quia minus mali puto esse quam in illa profectione, maneo in sententia et gaudeo te mansisse.

\ciceroSection{9} cum vero iam angerer et timerem ne quid a me dedecoris esset admissum, iii Nonas Mart., tamen te non esse una cum Pompeio non fero moleste. postea , si opus fuerit, non erit difficile, et illi, quoquo tempore fiet, erit ἀσμένιστον ; sed hoc ita dico si hic qua ratione initium fecit eadem cetera aget, sincere, temperate, prudenter, valde videro et consideratius utilitati nostrae consuluero.

\ciceroSection{10} vii Idus Martias scribis Peducaeo quoque nostro probari quod quierim, cuius auctoritas multum apud me valet. his ego tuis scriptis me consolor ut nihil a me adhuc delictum putem. tu modo auctoritatem tuam defendito; adversus me nihil opus est sed consciis egeo aliis. ego si nihil peccavi, reliqua tuebor. ad ea tute hortare et me omnino tua cogitatione adiuva. hic nihildum de reditu Caesaris audiebatur. ego his litteris hoc tamen profeci, perlegi omnis tuas et in eo acquievi.
